\scanpage{277}
\margin{\unsure{2}}
% archive foliation: 276

d'Orsay. Celle-ci est formée

a) des maîtres de conférences et professeurs des deux sections (Laboratoire de Math. et Ecole de Math.)

b) des assistants et maîtres assistants ou de leurs représentants élus, à raison d'un sixième au plus de l'Assemblée.

c) des représentants élus des étudiants, pour un tiers de l'Assemblée.

Les décisions sont prises à la majorité absolue des votes exprimés, sauf pour les nominations ou pour les changements organiques, pour lesquels une majorité des deux tiers des votes exprimés est requise.

4°/ a) \uline{Obligations} des membres du \uline{Laboratoire de Recherche}.

Les obligations d'enseignement de Maîtres de Conférences et Professeurs sont de trois heures hebdomadaires dont la moitié consiste obligatoirement en un cours de 3\textsuperscript{ème} cycle, l'autre moitié pouvant être

--- soit un autre cours de 3\textsuperscript{ème} cycle

--- soit la direction d'un séminaire

--- soit un enseignement dans le cadre de l'Institut d'Enseignement (de préférence dans la Maîtrise de Mathématiques)

En outre la recherche et la direction de recherches font partie des obligations des Maîtres de Conférences et Professeurs

b) \uline{Obligations} des membres de l'\uline{Ecole de Mathématique}.

Le service d'un Maître de Conférences ou Professeur de
